\documentclass[11pt]{beamer}

% ============================================================
% Theme
% ============================================================
\usetheme{Madrid}
\usecolortheme{seahorse}
\setbeamertemplate{navigation symbols}{}

% ============================================================
% Encoding (pdfLaTeX-safe)
% ============================================================
\usepackage[utf8]{inputenc}
\usepackage[T1]{fontenc}
\usepackage{lmodern}

% ============================================================
% Packages
% ============================================================
\usepackage{amsmath}
\usepackage{listings}
\usepackage{xcolor}

\lstset{
	basicstyle=\ttfamily\small,
	frame=single,
	breaklines=true,
	columns=fullflexible,
	keywordstyle=\color{blue},
	commentstyle=\color{gray},
	stringstyle=\color{red},
	showstringspaces=false
}

% ============================================================
% Title
% ============================================================
\title{Concurrency \& Distributed Systems}
\subtitle{Week 1 --- Part 5: Case Studies: Event Loops, Locks, Queues, and Contention}
\author{Dr. Mohamad Aoude}
\institute{Lebanese University \\ Faculty of Engineering}
\date{\today}

\begin{document}
	
	% ============================================================
	\begin{frame}
		\titlepage
	\end{frame}
	
	% ============================================================
	\section{Why Study These Architectures?}
	% ============================================================
	
	\begin{frame}{From Threads to Architecture}
		We move from:
		\begin{itemize}
			\item ``Add threads''
			\item to
			\item ``Choose the right concurrency model''
		\end{itemize}

		
		Four case studies:
		\begin{itemize}
			\item Single-thread event loop (Redis style)
			\item Lock-based shared state (contention)
			\item Queue growth (Little's Law)
			\item Read/write contention (SQLite model)
		\end{itemize}
	\end{frame}
	
	% ============================================================
	\section{Program 1: SingleThreadEventLoop (Redis style)}
	% ============================================================
	
	\begin{frame}{Program 1: Redis-Style Event Loop}
		\textbf{Idea:}
		\begin{itemize}
			\item One thread processes requests in a loop
			\item No shared-state races in request handling
			\item No locks around the ``work''
			\item Works well when per-request work is tiny (in-memory)
		\end{itemize}

		
		\begin{block}{Goal}
			Show that single-threaded can be \emph{fast} and \emph{predictable}.
		\end{block}
	\end{frame}
	
	\begin{frame}[fragile]{Core Idea (Simplified)}
		\begin{lstlisting}[language=Java]
	// Single-threaded event loop (Redis-style)
	while (running) {
		Request r = queue.poll();       // take next request
		if (r != null) {
			int x = 1 + 1;              // tiny in-memory work
			// respond
		}
	}
		\end{lstlisting}

		
		\begin{block}{Observation}
			No locks in the processing path. Predictable latency until overload.
		\end{block}
	\end{frame}
	
	\begin{frame}[fragile]{How to Run (Program 1)}
		\textbf{Compile \& run:}
		\begin{lstlisting}[language=bash]
			javac SingleThreadEventLoop.java
			java SingleThreadEventLoop
		\end{lstlisting}

		
		\textbf{What to observe:}
		\begin{itemize}
			\item high throughput
			\item stable latency when tasks are tiny
			\item failure mode is queue growth (not races)
		\end{itemize}
	\end{frame}
	
	% ============================================================
	\section{Program 2: ThreadedCounterWithLock}
	% ============================================================
	
	\begin{frame}{Program 2: Shared State + Lock Contention}
		\textbf{Idea:}
		\begin{itemize}
			\item Many threads update one shared counter
			\item Correctness requires a lock
			\item Lock serializes the critical section
		\end{itemize}

		
		\begin{alertblock}{Hypothesis to test}
			``More threads should make it faster.''
		\end{alertblock}
	\end{frame}
	
	\begin{frame}[fragile]{Critical Section}
		\begin{lstlisting}[language=Java]
	lock.lock();
	try {
		counter++;               // shared mutable state
	} finally {
		lock.unlock();
	}
		\end{lstlisting}

		
		\begin{alertblock}{Key point}
			All threads line up here. This becomes a bottleneck.
		\end{alertblock}
	\end{frame}
	
	\begin{frame}[fragile]{How to Run (Program 2)}
		\textbf{Compile:}
		\begin{lstlisting}[language=bash]
			javac ThreadedCounterWithLock.java
		\end{lstlisting}
		
		\pause
		
		\textbf{Run with increasing thread counts:}
		\begin{lstlisting}[language=bash]
			java ThreadedCounterWithLock 1
			java ThreadedCounterWithLock 2
			java ThreadedCounterWithLock 4
			java ThreadedCounterWithLock 8
			java ThreadedCounterWithLock 16
		\end{lstlisting}

		
		\textbf{Expected:}
		\begin{itemize}
			\item throughput rises a bit
			\item then plateaus
			\item then degrades (contention + context switching)
		\end{itemize}
	\end{frame}
	
	% ============================================================
	\section{Program 3: QueueingLittleLawDemo}
	% ============================================================
	
	\begin{frame}{Program 3: Little's Law and Queue Growth}
		Little's Law:
		\[
		L = \lambda W
		\]

		
		\begin{itemize}
			\item $L$: average number of requests inside the system
			\item $\lambda$: arrival rate (requests/sec)
			\item $W$: time in system (service + waiting)
		\end{itemize}

		
		\begin{alertblock}{Interpretation}
			If $\lambda$ rises (or $W$ rises), queues grow.
		\end{alertblock}
	\end{frame}
	
	\begin{frame}[fragile]{How to Run (Program 3)}
		\textbf{Compile \& run:}
		\begin{lstlisting}[language=bash]
			javac QueueingLittleLawDemo.java
			java QueueingLittleLawDemo
		\end{lstlisting}

		
		\textbf{Experiment: overload it}
		\begin{itemize}
			\item Open \texttt{QueueingLittleLawDemo.java}
			\item Increase \texttt{arrivalRate} (e.g., 50 $\rightarrow$ 200)
			\item Run again and observe \texttt{Remaining in queue}
		\end{itemize}
	\end{frame}
	
	\begin{frame}{What Students Must Conclude (Program 3)}
		\begin{itemize}
			\item If service time stays the same and arrivals increase, the queue must grow.
			\item Event loops (Redis style) \textbf{must keep work per request tiny}.
			\item Blocking I/O inside the event loop increases $W$ dramatically.
		\end{itemize}

		
		\begin{block}{Bridge}
			This is why Redis avoids blocking calls in the main loop.
		\end{block}
	\end{frame}
	
	% ============================================================
	\section{Program 4: SQLiteConcurrencyDemo}
	% ============================================================
	
	\begin{frame}{Program 4: SQLite Model (Many Readers, One Writer)}
		\textbf{SQLite-like principle:}
		\begin{itemize}
			\item Many readers can proceed together
			\item Writers serialize (exclusive access)
		\end{itemize}

		
		We simulate this with \texttt{ReentrantReadWriteLock}:
		\begin{itemize}
			\item read lock for readers
			\item write lock for writers
		\end{itemize}
	\end{frame}
	
	\begin{frame}[fragile]{Read/Write Lock Pattern}
		\begin{lstlisting}[language=Java]
	lock.readLock().lock();
	try {
		// read shared state
	} finally {
		lock.readLock().unlock();
	}
			
	lock.writeLock().lock();
	try {
		// write shared state
	} finally {
		lock.writeLock().unlock();
	}
		\end{lstlisting}
		
		\begin{alertblock}{Observation}
			Writers block everyone. Contention creates latency spikes.
		\end{alertblock}
	\end{frame}
	
	\begin{frame}[fragile]{How to Run (Program 4)}
		\textbf{Compile \& run:}
		\begin{lstlisting}[language=bash]
			javac SQLiteConcurrencyDemo.java
			java SQLiteConcurrencyDemo
		\end{lstlisting}

		
		\textbf{In class:}
		\begin{itemize}
			\item Increase number of writer threads in the code
			\item Observe how the system becomes ``spiky'' (more blocking)
		\end{itemize}
	\end{frame}
	
	% ============================================================
	\section{Grand Synthesis}
	% ============================================================
	
	\begin{frame}{Architectural Tradeoffs (The Point of the Case Study)}
		\begin{itemize}
			\item Single-thread event loop:
			\begin{itemize}
				\item simple, predictable
				\item fast for tiny in-memory tasks
				\item fails by queue growth under overload
			\end{itemize}
			\item Lock-heavy shared state:
			\begin{itemize}
				\item correctness cost (locks)
				\item contention and poor scaling
				\item higher variance / tail latency
			\end{itemize}
		\end{itemize}

		
		\begin{block}{Key insight}
			Engineering is tradeoffs. Single-threaded is a choice, not a limitation.
		\end{block}
	\end{frame}
	
	\begin{frame}{Student Questions (Submit as Short Answers)}
		Answer in 3--5 lines each:
		\begin{enumerate}
			\item Why is Redis single-threaded for request processing?
			\item When does a single-thread event loop fail?
			\item Why did throughput plateau in the lock-based counter?
			\item Use Little's Law: if $W$ doubles and $\lambda$ stays constant, what happens to $L$?
			\item Why does SQLite serialize writes?
		\end{enumerate}
	\end{frame}
	
	\begin{frame}
		\centering
		\Large Next Topic: Event-Driven I/O and Async Architectures
	\end{frame}
	
\end{document}
